\documentclass[12pt,letterpaper]{article}
\usepackage{amsmath,amssymb,amsthm}
\newtheorem{theorem}{Theorem}
\begin{document}
\title{Primes in Arithmetic Progressions}
\author{Hélène Dubois}
\date{November 2025}
\maketitle
\tableofcontents

\section{Introduction}
Euclid showed that there are infinitely many primes. Dirichlet proved much more: if $\gcd(a,q)=1$, then the progression $a, a+q, a+2q,\dots$ contains infinitely many primes.

\section{Statement}
\begin{theorem}[Dirichlet]\label{thm:dir}
Let $q\ge1$ and $\gcd(a,q)=1$. Then
\[
\sum_{p\equiv a\ (\mathrm{mod}\ q)}\frac1p=\infty.
\]
\end{theorem}

\subsection{Easy special cases}
The case $q=4$, $a=3$ has an elementary proof: if $p_1,\dots,p_k$ are primes of the form $4m+3$, the number $4p_1\cdots p_k-1$ has a prime factor of that form not among them.

\subsection{Quantitative form}
Let $\pi(x;q,a)$ count primes $p\le x$ with $p\equiv a\pmod q$. Then
\begin{equation}
\pi(x;q,a)\sim\frac{1}{\varphi(q)}\,\frac{x}{\log x}.
\end{equation}

\section{Remarks}
\paragraph{Characters.} The proof uses Dirichlet characters modulo $q$ and the nonvanishing of $L(1,\chi)$.
\paragraph{Computation.} A table of small cases for $q=10$ is easy to produce by hand.
\end{document}
